\section{Day 4: Expectation (Sep.\ 21, 2026)}
Consider a probability space $(\RR, \SB(\RR), \PP)$; then, given $X \colon \Omega \to \RR$ and $g \colon \RR \to \RR$, we have that
\[ \EE g(X) = \int_\Omega g(X(\omega)) \, d\PP(\omega) = \int_\RR g(x) \, d\PP_X(x). \]
We start with some simple examples.
\begin{enumerate}[(i)]
    \item Let $X \sim \opname{Ber}(p)$, i.e., $\PP(X = 1) = p$ and $\PP(X = 0) = 1 - p$. Then the expectation of $X$ can simply be computed as
    \[ \EE X = \int_\RR x \, d\PP_X(x) = 1 \PP_X(1) + 0 \PP_X(0) = p. \]
    In general, if $X$ is integer valued, then
    \[ \EE X = \sum_{n \in \ZZ} n \PP(X = n). \]

    \item Let $X$ be a standard normal random variable, i.e., $X$ has (Lebesgue) density
    \[ \rho(x) \coloneq \frac{e^{-\frac{x^2}{2}}}{\sqrt{2\pi}}. \]
    Then the expectation is as follows,
    \[ \EE X = \int_\RR x \rho(x) \, dx = 0, \]
    since $x \rho(x)$ is odd. We define $\EE X^k$ to be the $k^\text{th}$ moment of the random variable $X$. We now compute the variance.
    \[ \EE X^2 = \frac{1}{\sqrt{2\pi}} \int_\RR x^2 e^{-\frac{x^2}{2}} \, dx = \frac{-x e^{-\frac{x^2}{2}}}{\sqrt{2\pi}} \biggr\rvert_{-\infty}^\infty + \int_\RR \frac{e^{-\frac{x^2}{2}}}{\sqrt{2\pi}} \, dx. \]
    Since $\Var(X) \coloneq \EE X^2 - (\EE X)^2$ when $\EE X^2 < \infty$, we see $\Var(X) = 1$.
\end{enumerate}

\begin{fact}
    If $Y$ (a real random variable) has $\EE Y = a$, $\Var(Y) = b$, then $\EE(cY + d) = cd + d$, and $\Var(cY + d) = c^2 b$.
\end{fact}
Specifically, this is why we can assume, without loss of generality, that $\EE X = 0$ and $\Var(X) = 1$, since we may scale it however we want with these linear transformations. Specifically, if $X$ is a standard normal random variable, we say that $\sigma X + \mu$ is a normal random variable with mean $\mu$, variance $\sigma^2$, and write $\sigma X + \mu \sim N(\mu, \sigma^2$).

\begin{fact}
    If $X, Y$ are independent real random variables and $\EE \abs{X}, \EE \abs{Y} < \infty$, then $\EE \abs{XY} < \infty$ (i.e., $XY$ is integrable or $L^1$) and $\EE XY = \EE X \EE Y$.
\end{fact}
\begin{proof}
    We may directly compute
    \begin{align*}
        \EE \abs{XY} &= \int_{\RR^2} \abs{xy} d\PP_{(X, Y)}(x, y) \tag{$\PP(x, y)$ is the joint law of $(X, Y)$} \\
        &= \int_{\RR^2} \abs{xy} d\PP_X \times d\PP_Y(x, y) \tag{Independence} \\
        &= \left(\int_\RR \abs{x} d\PP_X\right) \left(\int_\RR \abs{y} d\PP_Y\right) \tag{Fubini and Tonelli} \\
        &= \EE\abs{X} \EE\abs{Y} < \infty.
    \end{align*}
    A similar proof shows that $\EE XY = \EE X \EE Y$.
\end{proof}

For further reading, see theorem 1.10 in Panchenko, i.e., the ``disintegration theorem'', for the non-product case.

Do infinite collections of random variables exist? As an example, suppose we want independent random variables $X_1, \dots, X_n$, each with CDF $F_i$. Then consider $(\Omega, \SA, \PP) = (\RR^n, \SB(\RR^n), \PP_{X_1} \times \dots \times \PP_{X_n})$, and take $X_i(\omega_1, \dots, \omega_n) = \omega_i$, i.e., the $i^\text{th}$ coordinate projection could take $\PP_{(x_1, \dots, x_n)}$, which would work for non-independent random variables.

\begin{theorem}[Kolmogorov consistency theorem]
    Let $T$ be a potentially uncountable index set. Consider the \textit{cylindrical $\sigma$-algebra} on $\RR^T$, i.e., the $\sigma$-algebra generated by all finite coordinate projections. Specifically, $C$ is generated by setes of the form $B \times \RR^{T \setminus N}$, where $B \in \SB(\RR^N)$, $\abs{N} < \infty$, and $N \subset T$.

    Suppose we have a family of probability measures $\{\PP^N, N \subset T, \abs{N} < \infty\}$ (on $\RR^N$), sayisfying a \textit{consistency condition}: for all $N \subset M \subset T$ with $\abs{N}, \abs{M} < \infty$ and all $B \in \SB(\RR^N)$, we have
    \[ \PP^M(B \times \RR^{M \setminus N}) = \PP^N(B). \]
    Then there exists a unique probability measure $\PP$ on $(\RR^T, C)$ so that $\PP(\omega \in \RR^T = \restr{\omega}{M} \in B) = \PP^M(B)$ for all $M \subset T$, $\abs{M} < \infty$, and $B \in \SB(\RR^M)$.
\end{theorem}

This will be proved in next lecture.